% ==============================================================================
% T0/FFGFT Theory: Document 242
% Title: The Scale Ladder – Building blocks and fields across all layers
%        of matter. A pedagogical sketch on inter-scale couplings in FFGFT.
% Author: Johann Pascher
% Date: May 2026
% Status: Hypothesis, not proof (plausibility sketch).
% References: Docs.~159 (concert hall / mode structure), 133 (K_frak and RG flow),
%             167 (LiNbO3), 168 (periodic system), 182 (universe maximal scale),
%             205 (narrative), 230 (Hilbert-space bridge), 241 (Two Layers)
% ==============================================================================

\section*{Preface}

This document is a \emph{plausibility sketch}, not a derivation. It
contains hypotheses, not proofs. Where a quantitative statement is
made, it is to be read as an order-of-magnitude estimate, not as a
prediction. The central question --- \emph{how strongly do the
different scales of matter couple geometrically with one another?}
--- is an open research question, not a solved problem.

The purpose of this document is not to deliver an answer, but to
\emph{state the question cleanly} --- so that anyone working through
FFGFT does not overlook the scale ladder between elementary particles
and cosmology as a gap, but recognizes it as research terrain.

FFGFT has so far been worked out \emph{at the lowest layer}:
elementary particles, modes on the 4D-torus, winding numbers. The
upper boundary --- the cosmological maximum scale --- is treated
separately in Doc.~182. What lies between --- atoms, molecules,
crystals, gases, liquids, planets, stars, galaxies --- is addressed
at scattered points in the document series (Doc.~167 for crystals,
Doc.~168 for atoms), but not developed narratively as a coherent
scale ladder. That is what this document tries to make up, in the
pedagogical style of Doc.~241.

\section{A generally unsolved problem}

Before FFGFT enters the picture, it is worth taking stock of what
standard physics actually says about this question --- or does not say.

The Standard Model (SM) of particle physics is extraordinarily precise
on its own scale. It describes quarks, leptons, and their interactions
with an accuracy no other physical theory matches. But it has
\emph{no hierarchy theory}: there is no derivation within the SM of
why quarks are bound in hadrons, hadrons in nuclei, nuclei in atoms
--- or why these structures look exactly as they do. That is an
observation, not an explanation.

\subsection*{What the Standard Model has}

Standard physics meets this problem through \emph{effective theories}:
QCD for quarks and gluons; nuclear physics as a phenomenological
extension; atomic physics via QED. The transitions between these
levels are not derived --- they are empirically adjusted. No one has
calculated from the QCD Lagrangian alone why a proton has the mass it
has, let alone why a hydrogen atom is built exactly as it is. Lattice
field theory can now simulate the proton numerically; \emph{deriving}
it --- understanding why exactly this structure and no other ---
remains out of reach.

The closest relative to a cross-scale explanation is the
\emph{renormalization group} (RG). The RG describes how coupling
constants scale with energy. Asymptotic freedom in QCD is precisely
that: at high energy quarks are quasi-free, at low energy strongly
bound. That is a genuine scaling effect. But the RG runs
\emph{vertically} through energy scales and says nothing about how the
structure at one scale \emph{shapes} the next. It describes the flow
of couplings, not the emergence of structures.

\subsection*{The actual gap}

The real question is: why does an atom, as a consequence of the quark
structure of the proton, carry exactly the structure it carries? The
SM cannot answer this. How the properties of the proton --- mass,
spin, magnetic moment --- emerge from QCD is still numerically hard
and conceptually open. And one level up: why do protons and neutrons
form exactly the nuclei they form, with exactly these binding
energies? Nuclear physics is empirical --- not derived from QCD.

This is not a failure of diligence. It is a structural feature of
standard physics: it describes the mechanisms \emph{within} each
level very well, and the \emph{local transitions} from one level to
the next through appropriately chosen coupling parameters. What it
does not describe is whether and how a geometry at one scale
\emph{propagates through} the scale ladder --- whether what is
fundamental at the quark level leaves any trace in the structure of
the atom or the crystal.

\subsection*{An empirical argument: external influence as structural proof}

There is an argument that requires no theory and yet points directly
to the heart of the question.

We influence molecular structures through external fields: an
electromagnetic field changes the conformation of a protein; a
mechanical pressure shifts a crystallisation pathway; a change in
temperature drives a chemical reaction in a different direction. This
influence is not a special effect --- it is experimental routine in
chemistry, biophysics, and materials science.

But it has a structural consequence that is easily overlooked:
\emph{for an external field at a higher scale to be able to alter a
molecular structure at a lower scale, coupling channels between these
scales must physically exist.} These channels are not a property of
the external field. They are a property of the structure itself.

And here comes the decisive step: coupling channels that physically
exist are bidirectional. If a macroscopic field configuration can act
on the electron shell of a molecule, then conversely the substructure
of the molecule --- the arrangement of atoms, the electron density,
the nuclear geometry --- acts on the superordinate structure. Not
through external intervention, but as system-internal propagation.

This is not a speculative conclusion. It already shows itself in
protein folding: a protein folds into its equilibrium structure
without external intervention. The information about \emph{how} it
folds resides in the amino-acid sequence --- at the chemical scale.
This information propagates upward to the macromolecular geometry.
That is inter-scale coupling, system-internal, without an external
source.

Standard physics describes these channels locally and correctly:
Coulomb interaction, van-der-Waals, hydrogen bonds. What it does not
describe is whether these local channels are the expression of a
\emph{continuous geometric structure} --- or merely a collection of
accidentally fitting mechanisms. FFGFT claims the former. The SM is
silent on the question.

\subsection*{Stability as a hidden presupposition}

Behind the SM's silence on inter-scale coupling lies a methodological
decision that is rarely made explicit: the SM treats stable
superordinate structures as \emph{boundary conditions}, not as
\emph{explanatory targets}. An atom exists --- fine, we calculate with
it. A crystal exists --- fine, we describe its band structure. Why
these structures are stable at all, and why they have exactly this
form and not another, is not asked. Stability is taken as given.

This conceals a fundamental assumption: that the inter-scale forces
which \emph{generate and maintain} this stability are somehow
self-evident. An atom is stable because the electromagnetic force
between nucleus and electrons has a certain strength and because
quantum mechanics stabilises certain orbitals. That is a precise
balance between forces at different scales --- nuclear forces at the
femtometre scale, electromagnetic at the ångström scale. This balance
is not trivial. It is the result of exactly that inter-scale coupling
which the SM does not explain. The SM \emph{uses} the stability that
emerges from this coupling without grounding it.

\subsection*{The planetary orbit argument}

In celestial mechanics it is taken for granted that a planetary orbit
is \emph{never} a pure two-body problem. Jupiter perturbs Saturn. The
Moon perturbs Earth's orbit. Tides couple rotation and orbital
evolution. Small couplings accumulate over millions of years into
structure-forming effects --- orbital resonances, the Laplace resonance
of Jupiter's moons, the stabilisation of Earth's axial tilt by the
Moon. No one in celestial mechanics says: ``We neglect Jupiter because
it is far away.'' Astronomy has learned that even small couplings over
long timescales are \emph{structure-constitutive}.

At the molecular and atomic level the same logic is not applied.
There, scale separation is the working principle --- QCD for quarks,
QED for atoms, solid-state physics for crystals, with carefully drawn
boundaries between these theories. The question of whether small
residual couplings between these scales become structurally visible
over sufficiently long timescales or in sufficiently precise
measurements is never even raised. Not because the answer is ``No'',
but because the effective theories have worked so well that no one has
asked what they silently integrate away.

This is not a physical necessity but a historically grown convention.
Celestial mechanics and particle physics treat inter-scale couplings
asymmetrically --- not from a principle, but from different traditions.
FFGFT claims to unify both under the same geometric logic: couplings
exist at all scales, their strength scales geometrically with the
logarithmic distance, and the stability of superordinate structures is
not self-evident but \emph{stands in need of explanation}.

\subsection*{Why astronomers see couplings that particle physicists do not}

There is a simple reason why inter-scale couplings are so conspicuous
in celestial mechanics and so invisible in particle physics --- and it
has nothing to do with whether the couplings exist in principle. It
has to do with \emph{how far apart} the scales involved actually are.

In celestial mechanics, Jupiter and Saturn play on the same stage.
They are both planets, both in the solar system, both on comparable
orbits. The distance between their scales is tiny --- almost zero when
you consider the full span of nature's scales. That is why their
mutual influence is large enough to leave visible traces over millions
of years.

In particle physics the scales between which one might look for
couplings are unimaginably far apart. Between a quark and an atom fit
eight powers of ten. Between the Planck scale and the scale of the
Higgs particle fit seventeen. A coupling that acts across so many
orders of magnitude grows weaker with every additional order ---
so weak that no instrument we have today could detect it.

This is the real reason why standard physics gets along so well with
separated scales: not because the couplings are absent, but because
the distances are so vast that the couplings practically vanish.
Astronomy and particle physics do not contradict each other --- they
are looking at different sections of the same ladder.

\subsection*{Simplicity at the bottom --- why fundamental constituents are so stable}

There is a second reason why the picture becomes ever more stable
going downward. It concerns not the distances between scales but the
nature of what sits at each level.

The deeper one descends on the scale ladder, the \emph{simpler} the
structures become. An electron has no inner parts, no substructure,
no degrees of freedom into which it could rearrange itself. A quark
is similarly simple. This simplicity is not accidental --- it is the
reason for their extraordinary stability. A simple structure cannot
decay because there is nothing simpler for it to decay into. The
lowest state at a given level is stable for the straightforward
reason that there is nothing below it.

The higher one climbs on the ladder, the more complex the structures
become. An atom already has many degrees of freedom --- electrons on
different orbitals, nuclear spin, vibrational modes. A molecule has
more. A protein can fold into thousands of different conformations.
At the planetary scale the complexity is so enormous that any
underlying geometric structure is buried under the noise of billions
of atoms and their interactions.

This growth of complexity upward explains why stability becomes
relative: an atom is stable under ordinary conditions but can be
ionised with enough energy. A molecule is stable at room temperature
but can be broken apart by heat or light. Stability diminishes the
further one moves from the simple building blocks at the bottom.

\subsection*{The inversion: how much effort it takes to break simple structures}

The same situation becomes particularly tangible when viewed from the
other direction.

To detach an electron from a hydrogen atom, ordinary ultraviolet light
suffices --- the kind of energy a lamp can produce. To split an atomic
nucleus requires neutron bombardment or extreme temperatures: millions
of degrees, such as those that prevail in stars or in a hydrogen bomb.
That is a million times more energy. To break a proton into its quarks
would require one of the largest particle accelerators in the world ---
and even then the quarks are never seen freely, because they
immediately form new bonds.

This staggering of effort is not accidental. It shows directly how
deep a structure sits on the scale ladder. The deeper, the simpler ---
and the more energy is needed to break it apart, precisely because
there is nothing it could ``fall'' into. The energy threshold is, in
a sense, the experimental fingerprint of depth.

Read in reverse: the fact that nuclear fission demands a million times
more effort than ionising an atom is direct proof that atomic nuclei
sit on a deeper, simpler level than electrons. And that even this is
not the bottom level is shown by the further effort needed to isolate
quarks. The scale ladder is not a metaphor --- it is physically
inscribed in these energy thresholds.

Standard physics describes these thresholds precisely and correctly.
What it does not explain is why the spacings between the scales are
exactly as large as they are --- why not half as large or twice as
large. FFGFT claims that these spacings follow a geometric order
rather than being accidental. Whether that is true is open. That the
question is legitimate, the staggering itself makes plain.

\subsection*{The hierarchy problem as symptom}

The most familiar symptom of this gap is the \emph{hierarchy problem}:
why is the gravitational scale (Planck mass) so many orders of
magnitude larger than the electroweak scale (Higgs mass)? The SM has
no explanation --- it adjusts the fine-tuning empirically to the
observed value. Supersymmetry, technicolour, Randall--Sundrum models:
all are attempts to explain this scale gap mechanically. None has been
experimentally confirmed so far.

The question this document asks --- how do the scales of matter couple
structurally to one another? --- is therefore not only a
FFGFT-specific gap. It is a \emph{generally unsolved question in
theoretical physics}. FFGFT names a mechanism intended to address it:
the continuous $\xi$-geometry on the 4D-torus, which carries the same
basic structure at every scale, with scale-specific corrections via
$K_{\mathrm{frak}}$. Whether this mechanism provides the answer is
open. That the question is open is not.

\section{The two reading perspectives}

If one wants to describe the world, there are two possible entry
perspectives. Both are old, both have a tradition in physics, and
both work in their own way.

\subsection*{Perspective 1 --- the building-block view}

The classical reductionist reading: the world consists of quarks;
quarks build nucleons; nucleons build atomic nuclei; nuclei and
electrons build atoms; atoms build molecules; molecules build
crystals or gases or liquids; these build larger structures ---
minerals, rocks, planetary shells, atmospheres, hydrospheres --- and
finally planets, stars, galaxies, galaxy clusters, the universe as a
whole. Each rung is a collection of the previous one. The binding
forces between the building blocks are described locally: the strong
interaction at the nuclear scale, the electromagnetic at the atomic
and molecular, van-der-Waals and hydrogen bonds at the
chemical-structural, gravity at the astronomical.

Standard physics moves largely in this perspective. It describes the
mechanisms \emph{within} each rung very well. It also describes the
\emph{transitions} between the rungs --- solid-state physics for how
atoms become crystals, hydrodynamics for how molecules become
liquids, astrophysics for how matter becomes stars. What it does
not describe is whether and how a geometry on one rung acts
\emph{through} the other rungs.

\subsection*{Perspective 2 --- the field view}

The holistic reading, central in FFGFT: all rungs are \emph{standing
waves} of the same underlying vacuum field. What appears as a
``building block'' on one rung is, at a deeper level, a stabilized
resonance configuration --- an electron is a mode on the 4D-torus, an
atom is a composite mode of electron and nuclear resonances, a
crystal is a periodic mode over atomic modes, a planet is a
gravitationally stabilized resonance of crystal aggregates. The same
geometry, the same winding logic, the same mode structure holds at
every rung --- only at a different scale.

The FFGFT claim here is not that the building-block perspective
would be wrong. It is not wrong. But it is a \emph{local language}
that has breaks between rungs which, in the field view, are no
breaks. The field view has a different strength: it shows why the
rungs lie \emph{just as} they lie --- not as a random accumulation
of building blocks, but as a self-consistent resonance chain.

\subsection*{Are both perspectives equivalent?}

Both describe the same world. They are not competitors, but two
languages of description for the same thing. Which is ``more
correct'' is a misguided question --- as little as a piano piece is
more correctly described by a score or by an acoustic frequency
spectrum. Both are exact; but they see different different things.

\section{The scale ladder}

If one orders the rungs --- from smallest to largest ---, there
emerges a logarithmic ladder whose rungs are roughly evenly spaced.
Each rung has a characteristic length, a characteristic binding
energy, and a logarithmic distance to the Planck scale.

\begin{center}
\small
\adjustbox{max width=\textwidth}{%
\begin{tabular}{lcc}
\toprule
\textbf{Layer} & \textbf{char.~length}~$\ell$ & $\log_{10}(\ell/\ell_{P})$ \\
\midrule
Planck scale         & $10^{-35}$~m & $0$ \\
Quark / lepton       & $10^{-18}$~m & $\sim 17$ \\
Atomic nucleus       & $10^{-15}$~m & $\sim 20$ \\
Atom                 & $10^{-10}$~m & $\sim 25$ \\
Molecule             & $10^{-9}$~m  & $\sim 26$ \\
Crystal (unit cell)  & $10^{-9}$~m  & $\sim 26$ \\
Microstructure (grain boundary) & $10^{-6}$~m & $\sim 29$ \\
Macroscopic matter   & $10^{-3}$~m  & $\sim 32$ \\
Planetary body       & $10^{7}$~m   & $\sim 42$ \\
Star                 & $10^{9}$~m   & $\sim 44$ \\
Solar system         & $10^{13}$~m  & $\sim 48$ \\
Galaxy               & $10^{21}$~m  & $\sim 56$ \\
Galaxy cluster       & $10^{23}$~m  & $\sim 58$ \\
Universe (observable) & $10^{26}$~m & $\sim 61$ \\
\bottomrule
\end{tabular}%
}
\end{center}

The rungs are not uniform --- but they span a logarithmic range of
about \emph{61 orders of magnitude}. That is the full span over
which the FFGFT geometry would have to act, if the field view is to
hold throughout.

\subsection*{What standard physics describes}

At each rung there is an established description:

\begin{itemize}\setlength\itemsep{2pt}
  \item \textbf{Quark/lepton:} Standard Model of particle physics
        (QCD, electroweak).
  \item \textbf{Atomic nucleus:} nuclear physics (shell model,
        mean-field, ab initio for light nuclei).
  \item \textbf{Atom:} quantum mechanics with Coulomb interaction;
        Hartree-Fock, DFT.
  \item \textbf{Molecule:} quantum chemistry; ab-initio methods up to
        about 100 atoms.
  \item \textbf{Crystal/solid:} solid-state physics; band model,
        Bloch theorem, collective excitations (phonons, magnons).
  \item \textbf{Liquid/gas:} statistical mechanics, hydrodynamics,
        thermodynamics.
  \item \textbf{Planetary body:} geophysics, mineralogy, classical
        mechanics.
  \item \textbf{Star/galaxy:} astrophysics, gravitational physics,
        cosmology.
\end{itemize}

Each of these theories is very accurate at its scale. But each works
with its \emph{own vocabulary}, its own constants, its own
interactions. Standard physics gives no continuous language
connecting all scales --- it welds them together by transition
theories.

\subsection*{What FFGFT would add}

FFGFT would say: at \emph{every} of these rungs the same vacuum
parameter $\xi_{0} = 4/30000$ acts. It acts everywhere with the same
local strength (this is the statement of Layer 1, Doc.~241). But the
\emph{cumulative} effect over the scale flow between the rungs is
different, depending on the logarithmic distance. This is the natural
generalization of what Doc.~133 makes explicit for the electroweak
scale:
\[
  K_{\mathrm{frak}}(\Lambda_1 \to \Lambda_2) \;=\; 1 - \delta\!\cdot\!\xi\cdot\ln(\Lambda_1/\Lambda_2),
\]
with $\delta$ a scale-specific coefficient that depends on the
geometric averaging over the rungs. For the jump
Planck~$\to$~electroweak this cumulated correction is about
$1{.}3\%$. For jumps at higher scales the corrections are
\emph{smaller}, not larger --- see the next section.

\section{The central hypothesis: small but continuous couplings}

Here comes the central point of this document, and it is honestly
flagged as \textbf{hypothesis, not proof}.

\subsection*{Statement}

The FFGFT geometry acts on all scales, but the direct geometric
couplings \emph{between} scales far apart from one another are
\emph{very small}. This is a necessary consequence of the
architecture:

\begin{enumerate}\setlength\itemsep{2pt}
  \item The parameter $\xi$ itself is small ($\sim 10^{-4}$).
  \item Coupling between scales enters only via the logarithmic
        distance --- not the linear distance.
  \item Even within the 17-order-of-magnitude jump (Planck $\to$
        electroweak), the cumulated correction is only $1{.}3\%$.
  \item Between scales only a few orders apart (e.g., atom $\to$
        crystal: $\sim 1$ order of magnitude), the direct geometric
        coupling is smaller by factors of $\sim 100$ --- that is,
        $\sim 0{.}01\%$ or below.
\end{enumerate}

Concretely this means: a geometry at the atomic scale ``shines
through'' into the macroscopic world only as a correction at the
order of $10^{-4}$ to $10^{-6}$ relative to the dominant interaction
of the respective scale. Standard physics mostly does not ``see''
these corrections, because they sit below its resolution and are
swept away by local renormalization or averaging procedures.

\subsection*{Why this justifies the building-block perspective}

Precisely because the inter-scale couplings are small, the
building-block perspective of standard physics is \emph{locally very
accurate}. Whoever does atomic physics can ignore that the same
geometry also acts in a galaxy cluster; the influence lies beyond any
spectroscopic resolution. Whoever does crystallography can ignore that
the same vacuum parameter also appears in electroweak theory. The
scale separation of standard physics is not wrong; it is a very good
approximation.

FFGFT does not contradict the scale separation. It only adds a
further statement: the separation is \emph{not absolute}. At
sufficiently high precision, traces of the continuous geometry should
appear --- as systematic deviations that cannot be explained by local
mechanisms.

\subsection*{Where such traces might be observable}

The following areas are the most plausible candidates where
inter-scale FFGFT couplings could become experimentally visible, if
the hypothesis is correct:

\begin{itemize}\setlength\itemsep{2pt}
  \item High-precision spectroscopy of atoms and molecules
        (anomalous shifts beyond the QED corrections, order $10^{-6}$
        relative).
  \item Anomalous corrections in solid-state band structures that
        cannot be explained by local many-body effects.
  \item Correlations between macroscopic constants (e.g., material
        densities, critical points, phase-transition energies) and
        atomic quantities, beyond known scaling laws.
  \item Astronomical ratios (planetary spacings, stellar
        mass-radius relations) that systematically differ slightly
        from purely classical-gravitational predictions.
\end{itemize}

All four points are formulated as research directions, not as
predictions. No one has so far systematically searched for such
couplings in the magnitude range $10^{-4}$ to $10^{-6}$, because no
established theory expects them.

\section{The multi-storey concert hall picture}

Whoever knows the concert-hall picture from Doc.~159 can see here a
structural extension: \emph{the concert hall has multiple storeys}.

On each storey there is an acoustic room of its own, with its own
resonances, its own reverberation time, its own characteristic
frequency. A tone played on the ground floor has its dominant
resonances on the ground floor. But it penetrates faintly, through
walls and ceilings, into the first storey and the second. What is a
full sound on the ground floor is in the second storey only a weak
sympathetic vibration --- but it is not zero. And conversely: a loud
tone on the top floor penetrates faintly downwards.

That is how to picture the scales: atomic geometry has its
resonances at the atomic rung, solid-state geometry at the
solid-state rung, planetary geometry at the planetary rung. But
each rung vibrates faintly in the others --- because they all stand
on the same vacuum, in the same fractal geometry. The direct
couplings are small because between widely separated rungs many
``walls and ceilings'' damp them. But they are not zero.

In this language, standard physics is the theory that describes
each storey \emph{by itself} very well. FFGFT is the theory that
asks how the building sounds as a whole --- and that claims that
there are quiet cross-couplings which a single storey cannot
explain.

\section{The two languages and their synthesis}

This allows the question stated at the outset --- how do
building-block view and field view connect without metaphor? --- to
be answered in the following form:

\begin{itemize}\setlength\itemsep{2pt}
  \item The \emph{building-block view} describes each rung
        \emph{locally} and its direct transitions to the next-higher
        or next-lower rung. It is the language of local mechanisms.
  \item The \emph{field view} describes the common vacuum on which
        all rungs stand, and the continuous geometry that is the
        same on all rungs. It is the language of global
        consistency.
  \item Both are exact at their own resolution level. The
        building-block view is more accurate at describing a single
        rung. The field view is more accurate at the question of
        why the rungs lie as they lie, and how they are
        self-consistently connected.
  \item The direct inter-scale couplings are \emph{small} --- so
        small that the building-block view practically always
        suffices. But they are not zero --- and at sufficiently
        high experimental precision they should become visible as
        systematic corrections.
\end{itemize}

\section{What we do not say}

We do \emph{not} say:

\begin{itemize}\setlength\itemsep{2pt}
  \item that standard physics is wrong;
  \item that all mechanisms between scales can be derived from FFGFT
        alone;
  \item that the multi-storey concert-hall picture is a complete
        theory;
  \item that the coupling strengths are quantitatively known from
        FFGFT --- they are \emph{not} derivable from the current
        documents, only plausible as orders of magnitude.
\end{itemize}

We do say, honestly as hypothesis:

\begin{itemize}\setlength\itemsep{2pt}
  \item that the geometry of the vacuum acts on all scales;
  \item that the scale separation of standard physics is a very good
        but \emph{not perfect} approximation;
  \item that the inter-scale couplings are small but in principle
        searchable in high-precision experiments;
  \item that this aspect of FFGFT is open research terrain.
\end{itemize}

\section{What would need to be done}

To turn this sketch into a testable theory, three steps would be
necessary, to be taken in future documents:

\begin{enumerate}\setlength\itemsep{2pt}
  \item \emph{A quantitative form of the coupling formula} --- how
        does the cumulated correction $K_{\mathrm{frak}}^{(n,m)}$
        between two scales $\Lambda_{n}$ and $\Lambda_{m}$ depend
        explicitly on the logarithmic distance? Doc.~133 contains
        the case Planck $\to$ electroweak; the generalization is
        outstanding.
  \item \emph{A selection of concrete observation sites} where the
        inter-scale coupling would be best testable, and a
        prediction of the expected effect size.
  \item \emph{A comparison with standard physics} in which it is
        honestly marked where FFGFT makes a prediction that goes
        beyond the Standard Model --- and where it merely provides
        an alternative language for the same fact.
\end{enumerate}

All three steps are non-trivial. They are what lies between a
sketch such as this document and a fully developed theory. This
document carries out none of these steps; it only marks them as
necessary.

\section*{Summary}

\begin{itemize}\setlength\itemsep{2pt}
  \item FFGFT is so far worked out at the lowest layer (elementary
        particles) and at the highest (cosmology); the intermediate
        layers --- atoms, molecules, crystals, matter aggregates,
        planets, stars, galaxies --- are only addressed in scattered
        places.
  \item The world can be read in two complementary languages:
        building-block view (local, standard physics) and field view
        (global, FFGFT). Both are exact; they see different
        different things.
  \item The FFGFT geometry, by assumption, acts on all scales with
        the same local strength. The cumulative effect over scale
        distances is small because the distances are large: traces
        of order $10^{-4}$ to $10^{-6}$ relative.
  \item This justifies the scale separation of standard physics as
        a very good approximation --- but not as a principled
        boundary.
  \item Status: \textbf{hypothesis, not proof}. Quantitative
        generalization of $K_{\mathrm{frak}}$ to arbitrary scale
        pairs and concrete predictions are outstanding.
\end{itemize}

\medskip
\noindent\textit{Status: May 2026. This document is methodologically
pedagogical and explicitly marked as a plausibility sketch. It
contains no independent quantitative predictions.}
